
\section*{Appendix: Two-Stream Instability}\label{app:A}

\subsection{Governing cold--fluid equations}

For each beam $j=\{+,-\}$ with equilibrium drift $v_{0j}=\pm u_0$:

\begin{align}
    \partial_t n_j + \partial_x (n_j v_j) &= 0, \label{eq:continuity}\\
    m_e\bigl(\partial_t v_j + v_j \partial_x v_j \bigr) &= -e E, \label{eq:momentum}\\
    \partial_x E &= -\frac{e}{\varepsilon_0}\bigl(n_+ + n_- \bigr). \label{eq:poisson}
\end{align}

\subsection{Linearisation}

Write $n_j = n_0 + n_{1j}$, $v_j = v_{0j} + v_{1j}$, $E = E_1$, retain only first--order terms and remove derivatives of zeroth--order terms.

% why remove the  \cancel{m_e v_{1j} \partial_x v_{0j}} term?
\begin{equation}
    \cancel{\partial_t n_{0j}} + \partial_t n_{1j} + n_{0j} \partial_x v_{1j} + n_{1j} \partial_x v_{0j} + \cancel{n_{0j} \partial_x v_{0j}} + \cancel{n_{1j} \partial_x v_{1j}} = 0 \label{eq:continuity}
\end{equation}
%
\begin{equation}
    \cancel{m_e \partial_t v_{0j}} + m_e \partial_t v_{1j} + \cancel{m_e v_{0j} \partial_x v_{0j}} + m_e v_{0j} \partial_x v_{1j} + \cancel{m_e v_{1j} \partial_x v_{0j}} + \cancel{m_e v_{1j} \partial_x v_{1j}} = \cancel{-e E_0} - e E_1, \label{eq:momentum}\\
\end{equation}
%
\begin{equation}
    \cancel{\partial_x E_0} + \partial_x E_1 = -\frac{e}{\varepsilon_0}\bigl(\cancel{n_{0+}} + \cancel{n_{0-}} \bigr) -\frac{e}{\varepsilon_0}\bigl(n_{1+} + n_{1-} \bigr)\label{eq:poisson}
\end{equation}

So we get 
\begin{align}
    \partial_t n_{1j} + n_{0j} \partial_x v_{1j} + n_{1j} \partial_x v_{0j} &= 0, \label{eq:continuity}\\
    %
    m_e \partial_t v_{1j} + m_e v_{0j} \partial_x v_{1j}  &= - e E_1, \label{eq:momentum}\\
    %
    \partial_x E_1 &= -\frac{e}{\varepsilon_0}\bigl(n_{1+} + n_{1-} \bigr)\label{eq:poisson}
\end{align}

Rearranging to have the time derivatives on the LHS
\begin{align}
    \partial_t n_{1j} &= -n_{0j} ~\partial_x v_{1j} - n_{1j} \partial_x v_{0j}, \label{eq:continuity}\\
    %
    \partial_t v_{1j} &= - e E_1/m_e - v_{0j} ~\partial_x v_{1j}, \label{eq:momentum}\\
    %
    \partial_x E_1 &= -\frac{e}{\varepsilon_0}\bigl(n_{1+} + n_{1-} \bigr)\label{eq:poisson}
\end{align}


Fourier transforming these equations in space ($\partial_x \rightarrow ik$) we get

\begin{align}
\partial_t n_{1j} &= -i k ~n_{0j} ~v_{1j} -i k ~n_{1j} ~v_{0j}\\
%
\partial_t v_{1j} &= - e E_1/m_e - i k ~v_{0j}~ v_{1j}\\
%
ik~ E_1 &= -\frac{e}{\varepsilon_0}\bigl(n_{1+} + n_{1-} \bigr)
\end{align}

Replacing the electric field back
\begin{align}
\partial_t n_{1j} &= -i k ~n_{0j} ~v_{1j} - i k ~v_{0j}~n_{1j} \\
%
 \partial_t v_{1j}  &= - i k ~v_{0j}~ v_{1j} - \frac{i~e^2}{k ~m_e\varepsilon_0}\bigl(n_{1+} + n_{1-} \bigr)
\end{align}

Defining $\alpha \equiv e^2 /(m ~\varepsilon_0 ~k^2)$ and setting $n_{0+}=n_{0-}=n_0$ and $v_{0+}=v_0=-v_{0-}$
\begin{align}
\partial_t n_{1+} &= -i k (+~v_0~n_{1+} + n_{0} ~v_{1+} ) \\
%
\partial_t n_{1-} &= -i k (-v_0 ~n_{1-} + n_{0} ~v_{1-} ) \\
%
\partial_t v_{1+}  &= - i k \left( +v_0~ v_{1+} + \alpha \left(n_{1+} + n_{1-} \right) \right) \\
%
\partial_t v_{1-}  &= - i k \left( -v_0~ v_{1-} + \alpha \left(n_{1+} + n_{1-} \right) \right)
\end{align}

This can be written in a Schrödinger equation form 
\begin{equation}
    \partial_t \begin{bmatrix}
        n_{1+}\\n_{1-}\\v_{1+}\\v_{1-}
    \end{bmatrix} = -ik \begin{bmatrix}
        v_0 & 0 & n_0 & 0\\
        0& -v_0 & 0 & n_0\\
        \alpha & \alpha & v_0 & 0\\
        \alpha & \alpha & 0 & -v_0
    \end{bmatrix} \begin{bmatrix}
        n_{1+}\\n_{1-}\\v_{1+}\\v_{1-}
    \end{bmatrix}, ~\partial_t \psi = -i H \psi
\end{equation}

------

Fourier transforming these equations in time ($\partial_t \rightarrow -i \omega$) we get
%Insert the plane--wave ansatz $e^{i(kx-\omega t)}$ to obtain, for each beam,
\begin{align}
    (-\omega + k v_{0j})\, n_{1j} + k n_0 v_{1j} &= 0, \label{eq:lin_cont}\\
    (-\omega + k v_{0j})\, v_{1j} &= -\,\frac{e}{m_e} E_1. \label{eq:lin_mom}
\end{align}

Solving \ref{eq:lin_mom} for $v_{1j}$ and back--substituting into \ref{eq:lin_cont} gives
    
\begin{align}
    v_{1j} &= -\,\frac{e k}{m_e(\omega - k v_{0j})}\,\phi_1,\\
    n_{1j} &= -\,\frac{e k^{2} n_0}{m_e(\omega - k v_{0j})^{2}}\,\phi_1,
\end{align}
with $E_1 = -i k \phi_1$.

\subsection{Dielectric function}

\begin{align}
    n_{1j} = -\frac{e  n_0}{\varepsilon_0 m_e} ~\dfrac{k^{2}}{(\omega - k v_{0j})^{2}}~\phi_1 \varepsilon_0 = -\omega_p^2 ~\dfrac{k^{2}}{(\omega - k v_{0j})^{2}} \dfrac{E_1}{-ik} \varepsilon_0 = \omega_p^2 ~\dfrac{k^{2}}{(\omega - k v_{0j})^{2}} \dfrac{e}{(-ik)^2} \sum_j n_{1j}
\end{align}

Insert $\sum_j n_{1j}$ into Poisson's equation \ref{eq:poisson} to obtain the cold--beam dielectric function

\begin{equation}
    \boxed{\;
    \varepsilon(\omega,k) \;=\;
    1 - \frac{\omega_p^{2}}{(\omega - k u_0)^{2}} - \frac{\omega_p^{2}}{(\omega + k u_0)^{2}} \;=\; 0
    \;} .
    \label{eq:eps}
\end{equation}

\subsection{Quartic form}

Clearing denominators in \ref{eq:eps} yields
\begin{equation}
    \boxed{\;
    \bigl(\omega^{2} - k^{2} u_{0}^{2}\bigr)^{2}
    = 2 \omega_{p}^{2}\,\bigl(\omega^{2} + k^{2} u_{0}^{2}\bigr)
    \;} .
    \label{eq:quartic}
\end{equation}

The symmetry of the two counter--propagating beams makes \ref{eq:quartic} an even function of $\omega$ and $k$.

\subsection{Roots: real and imaginary parts}

Let $x \equiv \omega^{2}$.  Equation~\ref{eq:quartic} is quadratic in $x$:
\[
x^{2} - 2(\omega_p^{2} + k^{2} u_{0}^{2})\,x
+ k^{2} u_{0}^{2}\,\bigl(k^{2} u_{0}^{2} - 2\omega_p^{2}\bigr) = 0 .
\]
Solving gives
\begin{equation}
    \boxed{%
    \omega^{2}
    = \omega_p^{2} + k^{2} u_{0}^{2}
    \;\pm\;
    \omega_p \sqrt{\omega_p^{2} + 4 k^{2} u_{0}^{2}}
    }\, .
    \label{eq:roots}
\end{equation}

\subsection*{Real branches}

The ``$+$'' choice in \ref{eq:roots} is always positive, giving two purely real modes (fast and slow).

\subsection*{Unstable branch}

With the ``$-$'' sign, the term under the outer square root can become negative.  Introducing the dimensionless variable
$\xi \equiv k u_0 / \omega_p$,
purely \emph{imaginary} solutions $\omega = i \gamma$ exist when $\xi^{2} < 2$, with growth rate

\begin{equation}
    \gamma(\xi)
    = \omega_p \sqrt{-\bigl(\xi^{2} + 1\bigr) + \sqrt{1 + 4\xi^{2}} } .
    \label{eq:gamma}
\end{equation}

\begin{itemize}
    \item \textbf{Threshold:} instability occurs for $k~ u_0 < \sqrt{2}\,\omega_p$.
    \item \textbf{Maximum growth:} differentiating \ref{eq:gamma} shows the peak at $\xi_{\max} = \sqrt{3}/2$, with
    $\gamma_{\max} = \tfrac12\,\omega_p$ ($= \Omega_p/\sqrt{2}$ in terms of the \emph{total} plasma frequency $\Omega_p = \sqrt{2}\,\omega_p$). Since this is independent of both $k$ and $u_0$, once the plasma frequency $\omega_p$ is chosen in the PIC/Vlasov simulation, the dominant growth rate will always be the same.
\end{itemize}

\subsection{Summary}

\begin{itemize}
    \item \emph{Dispersion relation:} equation~\ref{eq:quartic};
    \item \emph{Explicit roots:} equation~\ref{eq:roots} (real) and equation~\ref{eq:gamma} (imaginary / growth);
    \item \emph{Instability window:} $0 < k u_0 < \sqrt{2}\,\omega_p$;
    \item \emph{Peak growth rate:} $\gamma_{\max} = \omega_p/2$ at $k u_0 = (\sqrt{3}/2)\,\omega_p$.
\end{itemize}



\section*{Appendix: Weibel Instability}\label{app:B}

\subsection{Equilibrium and Governing Equations}

We consider a cold, non--relativistic, charge--neutral electron--positron plasma formed by two counter--streaming beams.
Species indices are $s\in\{e,+\}$, with equilibrium values
\[
n_s^{(0)}=\frac{n_0}{2}, \qquad 
\bm u_s^{(0)}=\pm v_0\,\hat{\bm z}, \qquad
\bm E^{(0)}=\bm B^{(0)}=\bm 0 .
\]

The cold two--fluid and Maxwell equations read
\begin{align}
\partial_t n_s + \nabla\!\cdot\!\bigl(n_s\bm u_s\bigr) &= 0, \\
\partial_t \bm u_s + (\bm u_s\!\cdot\!\nabla)\bm u_s &= \frac{q_s}{m}
\bigl(\bm E + \bm u_s \times \bm B\bigr), \\
\nabla\times\bm E &= -\partial_t \bm B, \qquad
\nabla\times\bm B = \mu_0 \bm J + \frac{1}{c^{2}}\partial_t \bm E .
\end{align}

\section{Linearisation and Fourier Analysis}

Perturb every quantity as $g = g^{(0)} + g^{(1)}$ and assume plane waves
$\,\propto \exp\!\bigl[i(kx-\omega t)\bigr]\,$ with wavevector $\bm k = k\hat{\bm x}$, i.e.\ \emph{perpendicular} to the drift direction.
Only $E^{(1)}_z$ and $B^{(1)}_y$ survive; Faraday gives $B^{(1)}_y = kE^{(1)}_z/\omega$.

\subsection{Single--Species Algebra}

\subsection*{Continuity}

For each species,
\[
n^{(1)}_s = \frac{k n_0}{2\omega}\,u^{(1)}_{sx}.
\]

\subsection*{Momentum}

Using $u^{(0)}_{sz}=\pm v_0$ and $q_e=-e,\;q_+=+e$,
\begin{align}
-i\omega\,u^{(1)}_{sx} &= \frac{q_s}{m}\,u^{(0)}_{sz}\,B^{(1)}_y, \\
-i\omega\,u^{(1)}_{sz} &= \frac{q_s}{m}\Bigl(E^{(1)}_z + u^{(1)}_{sx} B^{(1)}_y\Bigr).
\end{align}
Substituting $B^{(1)}_y$ and eliminating $u^{(1)}_{sx}$ yields
\[
u^{(1)}_{sz} = \frac{\mp e}{m\omega}\,E^{(1)}_z .
\]

\subsection{Current and Amp\`ere's Law}

The perturbed current along $z$ is
\[
J^{(1)}_z = \sum_s q_s n_0 u^{(1)}_{sz}
          = -\frac{e^{2} n_0}{m\omega}\,E^{(1)}_z .
\]

Using the curl--curl form of Maxwell's equations,
\[
k^{2} E^{(1)}_z = \frac{\omega^{2}}{c^{2}}\,E^{(1)}_z + \frac{i\omega\mu_{0}}{c^{2}}\,J^{(1)}_z .
\]
Insert $J^{(1)}_z$ and the definition $\omega_{p}^{2}\equiv 2e^{2}n_0/(\varepsilon_0 m)$ to obtain the \textbf{bi--quadratic}
\begin{equation}
\boxed{\;
\omega^{4} - \omega^{2}\!\bigl(\omega_{p}^{2}+k^{2}c^{2}\bigr)
          - \omega_{p}^{2}k^{2}v_{0}^{2}=0
\;}
\label{eq:quartic}
\end{equation}
which is the cold--fluid Weibel (filamentation) dispersion relation for an $e^\pm$ plasma.

\subsection{Roots and Growth Rate}

Equation~\ref{eq:quartic} gives
\begin{equation}
\omega^{2} = \frac{1}{2}\Bigl[\omega_{p}^{2}+k^{2}c^{2}
      \pm \sqrt{\bigl(\omega_{p}^{2}+k^{2}c^{2}\bigr)^{2}
      +4\omega_{p}^{2}k^{2}v_{0}^{2}}
      \Bigr].
\label{eq:roots}
\end{equation}

\begin{itemize}
  \item The ``$+$'' branch is always real and corresponds to a stable EM wave.
  \item The ``$-$'' branch becomes negative for suitable $k$, giving purely imaginary $\omega=i\gamma$ and hence exponential growth
  \[
  \gamma(k)=\frac{1}{\sqrt{2}}\,
  \sqrt{\sqrt{\bigl(\omega_{p}^{2}+k^{2}c^{2}\bigr)^{2}+4\omega_{p}^{2}k^{2}v_{0}^{2}}
  -\bigl(\omega_{p}^{2}+k^{2}c^{2}\bigr)}.
  \]
  In the long--wavelength limit $kc\ll\omega_{p}$ one recovers $\gamma\simeq k v_{0}$.
\end{itemize}

\subsection{Minimum Dimensionality for PIC Verification}

\begin{itemize}
  \item \textbf{Configuration space:} one spatial coordinate ($x$) suffices, giving a \emph{1D--3V} simulation.
  \item \textbf{Velocity space:} all three momentum components are required; a 2--V reduction would remove the Lorentz coupling driving the instability.
\end{itemize}

